% Two moves on one fixture tree, drawn by hand for Section~\ref{sec:phylo:search}.
% Input by textbook.tex; a sketch and not a rendered figure, so it is outside
% the figure cache and the render cap, which govern generated output only.
%
% Each panel is one unrooted topology on five taxa, drawn as three internal
% nodes on a spine with pendant edges. The macro takes the five leaf labels in
% the order (upper-left, lower-left, top-middle, upper-right, lower-right).
\providecommand{\moveleafpanel}[6]{%
  \begin{scope}[xshift=#1]
    \coordinate (u) at (0.0, 0); \coordinate (v) at (0.85, 0);
    \coordinate (w) at (1.7, 0);
    \draw[fglink] (u) -- (v) -- (w);
    \draw[fglink] (u) -- ++(-0.65, 0.6) node[fgnote, anchor=south east] {#2};
    \draw[fglink] (u) -- ++(-0.65,-0.6) node[fgnote, anchor=north east] {#3};
    \draw[fglink] (v) -- ++( 0.0, 0.8) node[fgnote, anchor=south] {#4};
    \draw[fglink] (w) -- ++( 0.65, 0.6) node[fgnote, anchor=south west] {#5};
    \draw[fglink] (w) -- ++( 0.65,-0.6) node[fgnote, anchor=north west] {#6};
  \end{scope}}
\begin{figure}[htbp]
\centering
\begin{tikzpicture}[scale=0.95]
  \node[font=\footnotesize, anchor=east] at (-1.0, 1.4) {NNI};
  \moveleafpanel{0cm}{A}{B}{C}{D}{E}
  \draw[->, thick] (2.85, 0) -- node[font=\footnotesize, above] {$e_1$} (3.75, 0);
  \moveleafpanel{4.3cm}{A}{C}{B}{D}{E}
  \draw[->, thick] (7.15, 0) -- node[font=\footnotesize, above] {$e_2$} (8.05, 0);
  \moveleafpanel{8.6cm}{A}{C}{D}{B}{E}
  \begin{scope}[yshift=-3.4cm]
    \node[font=\footnotesize, anchor=east] at (-1.0, 1.4) {SPR};
    \moveleafpanel{0cm}{A}{B}{C}{D}{E}
    \draw[->, thick] (2.85, 0) -- node[font=\footnotesize, above] {$B$} (3.75, 0);
    \moveleafpanel{4.3cm}{A}{C}{D}{B}{E}
    \draw[->, thick] (7.15, 0) -- node[font=\footnotesize, above] {$D$} (8.05, 0);
    \moveleafpanel{8.6cm}{A}{D}{C}{B}{E}
  \end{scope}
\end{tikzpicture}
\caption{Three successive states of one five-taxon topology under each move
set, from the same start. Above: two nearest-neighbour interchanges, each
swapping the subtrees adjacent to the named internal edge, so exactly one
bipartition changes per step. Below: two prune-and-regraft moves, each
detaching the named taxon and reattaching it at another edge. The first
regraft reaches in one move the topology the interchanges reach in two --- it
is the rightmost panel above --- which is what the containment
$\neighbourhood_{\mathrm{NNI}} \subset \neighbourhood_{\mathrm{SPR}}$ buys,
at a neighbourhood of $2(n-3)(2n-7) = 12$ candidates against $2(n-3) = 4$.}
\label{fig:nnispr}
\end{figure}
